\chapter{Friedmann Equation}

\section*{Introduction}

The Friedmann equation is $(\dot{a}/a)^2 = (8\pi G/3)\rho - k/a^2 + \Lambda/3$. Here $a(t)$ is the scale factor (normalized to $a = 1$ today), $\dot{a}/a = H$ is the Hubble parameter, $\rho$ is the total energy density (matter, radiation, and any other components), $k$ is the spatial curvature constant ($k = +1, 0, -1$ for closed, flat, and open universes), and $\Lambda$ is the cosmological constant (dark energy). The critical density $\rho_c = 3H^2/(8\pi G)$ determines the geometry: if $\rho = \rho_c$, the universe is flat ($k = 0$); if $\rho > \rho_c$, it is closed; if $\rho < \rho_c$, it is open. The density parameter $\Omega = \rho/\rho_c$ provides a dimensionless measure of how close the universe is to flat.
\section*{Fundamental Equation Used}

\[
\left(\frac{\dot{a}}{a}\right)^2 = \frac{8\pi G}{3}\rho - \frac{k}{a^2} + \frac{\Lambda}{3}
\]
\section*{Assumptions}

\textbf{Assumptions:} The universe is homogeneous and isotropic on large scales (the cosmological principle). General relativity is valid. The matter content can be described by a perfect fluid with equation of state $p = w\rho c^2$. The cosmological constant $\Lambda$ is truly constant in time.


\textbf{Limitations:} Does not explain the origin of the Big Bang (initial conditions must be supplied separately). Breaks down at the Planck scale ($t < 10^{-43}$ s), where quantum gravity effects dominate. Cannot describe anisotropic models (Bianchi universes) or inhomogeneities without perturbation theory. Assumes general relativity is the correct theory of gravity.


\section*{Mathematical Derivation}

\textbf{Step 1: The FLRW metric.}

The cosmological principle (homogeneity and isotropy on large scales) restricts the spacetime metric to the Friedmann--Lema\^etre--Robertson--Walker (FLRW) form:
\begin{align}
ds^2 = -c^2\,dt^2 + a(t)^2\left[\frac{dr^2}{1 - kr^2} + r^2(d\theta^2 + \sin^2\theta\,d\phi^2)\right],
\end{align}
where $a(t)$ is the scale factor (normalized so $a_0 = 1$ today) and $k = +1, 0, -1$ parametrizes spatial curvature. The Hubble parameter is $H = \dot{a}/a$.

\textbf{Step 2: Einstein field equations.}

The Einstein field equations are $G_{\mu\nu} + \Lambda g_{\mu\nu} = 8\pi G\,T_{\mu\nu}/c^2$. For a perfect fluid with energy density $\rho$ and pressure $p$, the stress-energy tensor is $T^{\mu\nu} = (\rho + p/c^2)u^\mu u^\nu + p\,g^{\mu\nu}$, where $u^\mu = (c, 0, 0, 0)$ in comoving coordinates.

\textbf{Step 3: Compute the 00-component.}

The time--time component of the Einstein tensor for the FLRW metric is
\begin{align}
G_{00} = \frac{3}{a^2}\left(\frac{\dot{a}^2}{c^2} + \frac{kc^2}{a^2}\right) \cdot c^2 = 3\left(\frac{\dot{a}^2}{a^2} + \frac{k}{a^2}\right).
\end{align}
Wait, let us be more careful. With $g_{00} = -1$, the 00-component gives:
\begin{align}
G^0_{\ 0} = -\frac{3}{c^2}\frac{\dot{a}^2 + kc^2/a^2}{a^2}\cdot a^2 \cdot c^2 = -3\left(\frac{\dot{a}}{a}\right)^2\frac{1}{c^2} - \frac{3k}{a^2}\cdot\frac{1}{c^2}\cdot c^2.
\end{align}
The $00$-component of the Einstein field equation $G_{\mu\nu} + \Lambda g_{\mu\nu} = 8\pi G\,T_{\mu\nu}/c^2$ gives:
\begin{align}
3\left(\frac{\dot{a}}{a}\right)^2 + \frac{3kc^2}{a^2} - \Lambda c^2 = \frac{8\pi G\rho c^2}{c^2} \cdot c^2 = 8\pi G\rho.
\end{align}
Wait, we should be more precise. Let us work in units where $c = 1$ and restore at the end. The 00-Einstein equation gives
\begin{align}
3\frac{\dot{a}^2 + k}{a^2} - \Lambda = 8\pi G\rho.
\end{align}
Rearranging:
\begin{align}
\left(\frac{\dot{a}}{a}\right)^2 = \frac{8\pi G}{3}\rho - \frac{k}{a^2} + \frac{\Lambda}{3}.
\end{align}

\textbf{Step 4: Physical interpretation of each term.}

The left side, $H^2 = (\dot{a}/a)^2$, is the squared expansion rate. On the right: $8\pi G\rho/3$ is the gravitational source term---matter and energy drive expansion or contraction. The curvature term $-k/a^2$ determines the spatial geometry: $k = 0$ gives flat (Euclidean) space; $k = +1$ gives a closed (spherical) universe; $k = -1$ gives an open (hyperbolic) universe. The cosmological constant term $\Lambda/3$ drives accelerated expansion if $\Lambda > 0$.

\textbf{Step 5: Critical density and the density parameter.}

Setting $\dot{a} = 0$ (a momentarily static universe) and $\Lambda = 0$ gives the critical density:
\begin{align}
\rho_c = \frac{3H^2}{8\pi G}.
\end{align}
Defining the density parameter $\Omega = \rho/\rho_c$, we can rewrite the Friedmann equation as
\begin{align}
\frac{k}{a^2 H^2} = \Omega - 1.
\end{align}
Thus $\Omega = 1$ implies $k = 0$ (flat), $\Omega > 1$ implies $k = +1$ (closed), and $\Omega < 1$ implies $k = -1$ (open). Current observations give $\Omega_{\text{total}} = 1.0007 \pm 0.0019$, consistent with a flat universe.

\section*{Characteristics of the Equation}

The Friedmann equation is a first-order ODE for $a(t)$. Its behavior depends on the relative contributions of the three terms. In a matter-dominated universe ($k = 0$, $\Lambda = 0$), $a(t) \propto t^{2/3}$. In a radiation-dominated universe, $a(t) \propto t^{1/2}$. In a $\Lambda$-dominated universe, $a(t) \propto e^{Ht}$ (exponential expansion). The transition from decelerating expansion (matter/radiation domination) to accelerating expansion ($\Lambda$ domination) occurred about 5 billion years ago. The Hubble parameter $H = \dot{a}/a$ is not constant; it evolves as the energy density changes.
\section*{Solution Methods}

\textbf{Analytical solutions:} For simple cases (single component with constant $w$), the Friedmann equation can be integrated exactly: $a(t) \propto t^{2/(3(1+w))}$. \textbf{Numerical integration:} For realistic cosmologies with multiple components (radiation + matter + $\Lambda$), integrate the Friedmann equation numerically using standard ODE solvers. \textbf{Parameter fitting:} Adjust $\Omega_m$, $\Omega_r$, $\Omega_\Lambda$, and $H_0$ to match observational data (CMB power spectrum, supernova distances, baryon acoustic oscillations). \textbf{Perturbation theory:} For structure formation, linearize the Friedmann equation around the homogeneous solution and solve for density perturbations.
\section*{Verifications}

Check the matter-dominated flat universe: $a(t) \propto t^{2/3}$ gives $H = \dot{a}/a = 2/(3t)$, so the age of the universe is $t_0 = 2/(3H_0)$. For $H_0 = 70$ km/s/Mpc, this gives $t_0 \approx 9.3$ Gyr (without $\Lambda$). With $\Lambda$, the age is older ($\sim 13.8$ Gyr) because the expansion was slower in the past. Verify that the critical density gives $k = 0$: $\rho_c = 3H_0^2/(8\pi G) \approx 9.5 \times 10^{-27}$ kg/m$^3$. Check that $\Omega_m + \Omega_\Lambda + \Omega_r = 1$ for a flat universe (consistent with CMB observations).
\section*{Applications}

The Friedmann equation is the foundation of modern cosmology. It predicts the expansion history of the universe, which is tested by measuring the distances to Type Ia supernovae (the discovery that led to the 2011 Nobel Prize). It determines the age of the universe ($\sim 13.8$ billion years). It predicts the cosmic microwave background (CMB) anisotropy spectrum, which has been measured to high precision by WMAP and Planck. It governs nucleosynthesis in the early universe, determining the primordial abundances of hydrogen, helium, and lithium. It also predicts the growth of large-scale structure (galaxies, clusters, voids) through gravitational instability.
\section*{What Richard Feynman Would Have Asked}

\subsection*{1. Plain-English Translation}
\textit{``What is the Friedmann equation actually telling us about the universe?''}

The Friedmann equation is a bookkeeping statement for the expansion of the universe. It says: the square of the expansion rate equals the sum of three competing effects---the gravitational pull of matter (which slows expansion), the curvature of space (which can either help or hinder expansion), and dark energy (which can drive acceleration). At any moment, the expansion rate is determined by the balance among these three. If matter dominates, the universe decelerates. If dark energy dominates, it accelerates.

Think of it as a cosmic tug-of-war. Matter pulls inward, trying to slow or reverse the expansion. Dark energy pushes outward, driving acceleration. Curvature determines the geometry of the arena where this contest plays out. The Friedmann equation tells you who is winning at any given epoch.

\subsection*{2. Localized Mechanics}
\textit{``If I look at a single galaxy cluster, what does the Friedmann equation tell me about its motion?''}

Locally, the expansion of the universe is irrelevant: galaxy clusters are bound by their own gravity and do not participate in the Hubble flow. The Friedmann equation describes the average expansion of the universe on scales of hundreds of megaparsecs, where the cosmological principle (homogeneity and isotropy) holds. On smaller scales, local gravitational dynamics (described by the full Einstein field equations or Newtonian gravity) dominate.

However, the Friedmann equation determines the background through which light from distant galaxies travels. The redshift of a galaxy's light is determined by how much the scale factor has increased since the light was emitted: $1 + z = a(t_{\text{emission}})/a(t_{\text{observation}})$. This is the cosmological redshift, and it is the primary observational evidence for expansion.

\subsection*{3. Testing Extreme Limits}
\textit{``What happens at the Big Bang, at infinite time, and for extreme values of $\Lambda$?''}

The Big Bang ($t \to 0$): The Friedmann equation gives $a(t) \to 0$, $\rho \to \infty$, and $T \to \infty$. The equation itself breaks down at the Planck time ($t \sim 10^{-43}$ s), where quantum gravity effects become dominant. We do not know what ``happened'' at or before the Planck time.

Infinite Time ($t \to \infty$): For $\Lambda > 0$, the universe expands exponentially: $a(t) \propto e^{Ht}$. Distant galaxies recede faster than the speed of light (their recession velocity is $v = H d$, which exceeds $c$ for $d > c/H$). They disappear beyond the cosmological horizon, and the observable universe becomes increasingly empty.

Extreme $\Lambda$: If $\Lambda$ is very large, the universe inflates exponentially even at early times (inflation). If $\Lambda$ is negative, the universe recollapses in a ``Big Crunch.'' The sign and magnitude of $\Lambda$ determine the qualitative fate of the universe.

\subsection*{4. Dimensionality and Geometry}
\textit{``Does the curvature parameter $k$ really matter? Isn't the universe flat?''}

The curvature parameter $k$ determines the global geometry of space: positive curvature ($k = +1$) gives a closed universe (finite volume, like the surface of a sphere); zero curvature ($k = 0$) gives a flat universe (infinite volume, Euclidean geometry); negative curvature ($k = -1$) gives an open universe (infinite volume, hyperbolic geometry). Observations (CMB, BAO) indicate that $k$ is very close to zero: $\Omega_k = -0.002 \pm 0.003$.

Even if $k$ is exactly zero today, it could have been non-zero in the early universe. Inflationary cosmology predicts that $k$ is driven exponentially close to zero by the rapid expansion in the first $10^{-32}$ seconds. The fact that $k \approx 0$ is one of the key pieces of evidence for inflation.

\subsection*{5. Cross-Disciplinary Connection}
\textit{``Where else does this equation show up outside of cosmology?''}

The Friedmann equation is a special case of the Raychaudhuri equation, which describes the convergence or divergence of geodesic congruences in general relativity. The same mathematical structure---tracking how a ``scale factor'' evolves under the influence of energy content and curvature---appears in the study of gravitational collapse (Oppenheimer--Snyder model), in the physics of the early universe (inflation, nucleosynthesis), and in analogue gravity models (sound waves in expanding fluids).

In fluid dynamics, the continuity equation and Euler equation for a homogeneous, isotropic, expanding fluid are mathematically identical to the Friedmann equation. This is not a coincidence---Friedmann's original derivation treated the universe as a perfect fluid. The same mathematics describes any system where a uniform medium expands or contracts under its own gravity.

%=============================================================
\section*{FAQ}

\textbf{Q: Does the universe expand ``into'' something?}

A: No. The Friedmann equation describes the expansion of space itself, not motion through space. The scale factor $a(t)$ multiplies all distances simultaneously. There is no ``outside'' that the universe is expanding into. The question ``what is the universe expanding into?'' is like asking ``what is north of the North Pole?''---it is a category error.

\textbf{Q: What is dark energy?}

A: Dark energy is the name given to whatever causes the accelerated expansion of the universe. The simplest candidate is the cosmological constant $\Lambda$, which corresponds to a constant energy density of the vacuum. However, the observed value of $\Lambda$ is $10^{120}$ times smaller than quantum field theory predicts---the ``cosmological constant problem,'' one of the deepest unsolved problems in physics.

\textbf{Q: Will the universe expand forever?}

A: With the observed values ($\Omega_\Lambda \approx 0.7$, $\Omega_m \approx 0.3$, $k \approx 0$), the universe will expand forever at an accelerating rate. The cosmological constant dominates, and $a(t) \propto e^{Ht}$ asymptotically. If dark energy evolves differently (phantom energy, $w < -1$), the expansion could lead to a ``Big Rip.''
\section*{Important Bibliography}

\begin{enumerate}
\item Weinberg, S., \textit{Gravitation and Cosmology: Principles and Applications of the General Theory of Relativity} (Wiley, 1972), Chapter 15.
\item Carroll, S.M., \textit{Spacetime and Geometry: An Introduction to General Relativity}, 3rd ed. (Cambridge University Press, 2024), Chapter 8.
\item Friedmann, A., ``\"Uber die M\"oglichkeit einer Welt mit konstanter negativer Kr\"ummung des Raumes,'' \textit{Zeitschrift f\"ur Physik} \textbf{21}, 326--332 (1924).
\item Planck Collaboration, ``Planck 2018 results. X. Constraints on inflation,'' \textit{Astronomy \& Astrophysics} \textbf{641}, A10 (2020).
\end{enumerate}


